% Candidate A.56. Notes only pending independent audits.
\subsection{A reduction criterion for arbitrary finite covers}

\begin{statement}{A.56}{Proved}{Cuspidal normalizers from special rank and at all sufficiently large Gaussian split primes}
Let \(Y\) be an affine hyperbolic curve over a number field,
with smooth proper compactification \(C\) of genus \(g\ge2\),
and let \(p\) be a prime. Suppose that \(C\) has potentially
good \emph{curve} reduction at some place above \(p\), and that
the good special Jacobian has \(p\)-rank at least two.
Then every nontrivial closed procyclic subgroup of
\(\pi_1(Y_{\overline K})^{(p)}\) with open arithmetic normalizer
is contained in cusp inertia.

More specifically, let \(Y\to U=E\setminus\{O\}\) be any
geometrically connected finite etale cover, where
\(E:y^2=x^3-x\), and let \(C\to E\) be its proper
compactification. If \(p\equiv1\pmod4\) and \(C\) has
potentially good curve reduction at a place above \(p\),
the same normalizer conclusion holds, with no special-rank
hypothesis on \(C\).

In particular, for every such cover of degree \(d\), its
pro-\(p\) normalizer criterion holds for every
\[
                         p>d,\qquad p\equiv1\pmod4.
                                                        \tag{A.56.1}
\]
These assertions concern actual normalized procyclic subgroups.
They do not assert full fixed-subgroup comparisons or a
cofinal family of covers.
\end{statement}
\begin{proof}
\paragraph{Choosing the Picard torsion without a CM splitting.}
After a finite number-field extension take a good proper curve
model at the selected place, and let \(J\) be its Jacobian.
Put
\[
 T=T_pJ,\qquad
 s:T\longrightarrow T_p^{\rm et}\overline J,
 \qquad K=\ker s.                                    \tag{A.56.2}
\]
The specialization map is surjective and its kernel is saturated.
Indeed the connected--etale sequence of the \(p\)-divisible group
over the complete strictly henselian valuation ring gives this
sequence on geometric generic Tate modules; see
\cite[Sections 1.4 and 2.2, printed pp.160 and164,
and Section 2.4, p.169]{fghTate1967Pdivisible}.
Surjectivity at each geometric generic finite level and the
surjective transition maps on the connected kernels give
surjectivity on their inverse limits. The etale quotient
specializes isomorphically to the geometric special Tate module.
If the special \(p\)-rank is \(f\), then
\(\operatorname{rank}T=2g\) and
\(\operatorname{rank}K=2g-f\). We do not assert that this
sequence splits as integral arithmetic representations.

Suppose \(A\) is a noncuspidal procyclic subgroup with open
arithmetic normalizer, choose a generator \(u\), and let
\(v=p^r v_0\ne0\) be its proper integral homology class,
with \(v_0\) primitive. Its existence and the stability of
its rational span follow from A.11. Use the principal
polarization to regard \(T\) as the symplectic integral
homology lattice. In \(T/pT\), the hyperplane
\(\overline v_0^\perp\) has dimension \(2g-1\).
Its intersections with \(K/pK\) and with the line
\(\F_p\overline v_0\) are both proper subspaces: for the
first use \(f\ge2\), and for the second use \(g\ge2\).
The union of two proper linear subspaces cannot be the whole
space, even over \(\F_2\). Choose
\[
 \overline w\in\overline v_0^\perp,
 \qquad \overline w\notin K/pK,
 \qquad \overline w\notin\F_p\overline v_0.             \tag{A.56.3}
\]
Since \(v_0\) is primitive and the symplectic pairing perfect,
the integral map \(T\to\Z_p(1)\) given by pairing with
\(v_0\) is surjective. Its kernel reduces onto the displayed
hyperplane. Thus \(\overline w\) lifts to
\(w\in v_0^\perp\), and \(s(w)\) is primitive.

For \(n>r\), let \(L_n\in J[p^n]\) be the point represented
by \(w\bmod p^n\). Its specialization has exact order
\(p^n\). Each finite-level point is defined over a finite
number field. After rigidification at a rational point it gives
an actual degree-zero line bundle on the good curve model,
with its \(p^n\)-th power trivialized, exactly as in A.55.
The resulting character \(\eta_n\), defined by Kummer
duality and polarization, satisfies
\[
       \eta_n(v)=0,
       \qquad v\bmod p^n\notin
               (\Z/p^n)\operatorname{PD}(\eta_n).     \tag{A.56.4}
\]
For the second assertion, an equality
\(p^r v_0=a w\pmod {p^n}\) forces \(p^r\mid a\),
because \(w\) is primitive. Divide by \(p^r\) and reduce
modulo \(p\), contradicting (A.56.3).

\paragraph{The actual cover and its new homology.}
The \(\mu_{p^n}\)-torsor defined by \(L_n\) gives an
unramified proper cyclic cover \(q_n:C_n\to C\). The
normality and radicial-main-component proof in A.55 applies
without change for every genus. The main component is
\(\overline C^{(p^{-n})}\); the norm on its Jacobian is a
Frobenius isogeny. Raynaud's Theorem 1 therefore identifies
the special Prym, up to isogeny, with the product of
exceptional component Jacobians, all with slopes strictly
between zero and one. By A.46 its Tate representation has
no rational line stable under an open Galois subgroup.

The subgroup \(A\) lies in the index-\(p^n\) covering
subgroup because \(\eta_n(v)=0\). Pass to the actual
punctured curve \(q_n^{-1}Y\), and intersect the normalizer
with its arithmetic fundamental group after a finite
descent extension. The arithmetic image remains open.
The class \(\widetilde v\) of the same generator \(u\)
in \(H_1(C_n,\Z_p)\) spans a stable line and pushes
forward to \(v\). Its Prym projection vanishes, so it is
in the old rational homology and is integrally invariant
under the cyclic deck group.

The integral calculation (A.55.2) works in every genus.
For a symplectic basis \(a_1,b_1,\ldots,a_g,b_g\)
adapted to \(\eta_n(b_1)=1\), the Schreier homology has
the two fixed generators \(h,B\), and two regularly
permuted \(p^n\)-tuples for each \(j\ge2\). Thus
\[
 q_{n,*}H_1(C_n,\Z_p)^{\operatorname{Deck}(q_n)}
 =\Z_p a_1\oplus p^n\Z_p b_1
       \oplus\bigoplus_{j=2}^g
             (p^n\Z_p a_j\oplus p^n\Z_p b_j).         \tag{A.56.5}
\]
This contradicts (A.56.4), proving the first assertion.

\paragraph{The once-punctured Gaussian elliptic target.}
Now let \(q:C\to E\) compactify the given cover. Its
nonconstant map makes \(E\) an isogeny factor of \(J_C\):
\(q_*q^*=[\deg q]\), so
\[
                     J_C\sim E\times P_q.              \tag{A.56.6}
\]
At \(p\equiv1\pmod4\), the good elliptic reduction of
\(E\) is ordinary. One can check this directly from the
nonzero central coefficient of \((x^3-x)^{(p-1)/2}\);
equivalently it is the usual split-prime ordinary
criterion for CM by \(\Z[i]\).
After a finite extension giving the good curve model of
\(C\), all factors in (A.56.6) have good reduction.
Special \(p\)-rank is additive in products and invariant
under isogeny, hence
\(f(J_C)=1+f(P_q)\ge1\).

If \(f(J_C)\ge2\), the first assertion applies.
If \(f(J_C)=1\), the Prym has special \(p\)-rank zero.
Statement A.46 forces any stable rational line in
\(V_pJ_C\) to project trivially to \(V_pP_q\).
It therefore projects nontrivially under \(q_*\) to
\(V_pE\). The actual punctured morphism
\(Y\to E\setminus\{O\}\) carries a proposed normalized
procyclic subgroup with this line to a nontrivial
arithmetically normalized subgroup with nonzero proper
homology. The already proved original elliptic criterion
puts its image in cusp inertia, whose proper homology
vanishes, a contradiction. This also covers proper
genus one, when \(q_*\) is an isogeny. Thus no genus
restriction is needed for the second assertion.

\paragraph{A degree bound which ensures good curve reduction.}
Suppose the original connected cover has degree \(d\),
and \(p>d\). The geometric Galois closure has group a
subgroup of \(S_d\), so its order is prime to \(p\).
The marked elliptic curve \((E,O)\) has good reduction
at this odd prime. The prime-to-\(p\) specialization
theorem for a smooth proper curve with a relative
normal-crossing boundary identifies its geometric
generic prime-to-\(p\) covers with the corresponding
special covers, preserving inertia; see
\cite[Expos\'e XIII, Corollary 2.9, Section 2.10,
and the proof of Corollary 2.12, new pp.288--292]{fghSGA1}.
Consequently, after a finite extension, the actual
Galois closure extends tamely over the good pair
\((\mathcal E,\mathcal O)\).

This extension has smooth proper source. Away from the
marked section it is etale. At a point above the section,
after an etale base change its completed local map is
\(R[[t]]\to R[[z]],\ t=z^e\), with \(e\) prime to
\(p\); the source is smooth over \(R\). The same
local description, or the subgroup quotient of this
tame cover, gives a smooth model for the original
non-Galois curve \(C\). Its good reduction is thus
potential, with no additional unspecified exceptional
prime set. Applying the second assertion gives
(A.56.1).
\end{proof}
